% FAC-EN-U0073
% Source: fac_body.tex lines 2959--3012.
% Authority: PDF physical pp. 74--75; printed pp. 269--270
% (Chapter III no. 75 complete; hard stop before no. 76).
% French control: no new transcription repair or printed-source correction.

\medskip
\phantomsection\label{fac:III:5:no:75}
\textbf{75. Nonsingular varieties.} The
\hyperref[fac:III:5:no:75:thm:3]{following result} plays an
essential role in extending the ``duality theorem'' of [15] to the
abstract case:

\phantomsection\label{fac:III:5:no:75:thm:3}
\textsc{Theorem 3.} \textit{Let $V$ be a nonsingular subvariety of the
projective space $P_r(K)$; suppose that all the irreducible components
of $V$ have the same dimension $p$. Let $\mathfrak F$ be a coherent
algebraic sheaf on $V$ such that, for every $x\in V$, $\mathfrak F_x$
is a free module over $\mathcal O_{x,V}$. Then
$H^q(V,\mathfrak F(-n))=0$ for all sufficiently large $n$ and
$0\leq q<p$.}

\phantomsection\label{fac:III:5:no:75:thm:3:proof}
By \hyperref[fac:III:5:no:74:thm:2]{Theorem~2}, it suffices to show that
$\mathcal O_{x,V}$, regarded as an $\mathcal O_x$-module, has dimension
$\leq r-p$. Let $\mathcal I_x(V)$ denote the kernel of the canonical
homomorphism
$\varepsilon_x:\mathcal O_x\longrightarrow\mathcal O_{x,V}$.
Since $x$ is a simple point of $V$, this ideal is known (cf. [18],
Theorem~1) to be generated by $r-p$ elements
$f_1,\ldots,f_{r-p}$; and the Cohen--Macaulay theorem
(cf. [13], p.~53, Proposition~2) shows that
\phantomsection\label{fac:III:5:no:75:formula:colon}
\[
(f_1,\ldots,f_{i-1}):f_i=(f_1,\ldots,f_{i-1})
\qquad\text{for }1\leq i\leq r-p.
\]

Let $L_q$ denote the free $\mathcal O_x$-module with basis elements
$e\langle i_1\cdots i_q\rangle$ corresponding to the sequences
$(i_1,\ldots,i_q)$ such that
\phantomsection\label{fac:III:5:no:75:def:Lq}
\[
1\leq i_1<i_2<\cdots<i_q\leq r-p.
\]
For $q=0$, take $L_0=\mathcal O_x$, and set
\phantomsection\label{fac:III:5:no:75:formula:differential}
\[
d\bigl(e\langle i_1\cdots i_q\rangle\bigr)
=\sum_{j=1}^{j=q}(-1)^j f_{i_j}\,
e\langle i_1\cdots\widehat{\imath}_j\cdots i_q\rangle
\]
and
\[
d\bigl(e\langle i\rangle\bigr)=f_i.
\]

By [6], Chapter~VIII, Proposition~4.3, the sequence
\phantomsection\label{fac:III:5:no:75:formula:resolution}
\[
0\longrightarrow L_{r-p}\overset d\longrightarrow L_{r-p-1}
\overset d\longrightarrow\cdots\overset d\longrightarrow L_0
\overset{\varepsilon_x}\longrightarrow\mathcal O_{x,V}
\longrightarrow0
\]
is exact. This proves that
$\dim_{\mathcal O_x}(\mathcal O_{x,V})\leq r-p$, as required.

\phantomsection\label{fac:III:5:no:75:cor}
\textsc{Corollary.} \textit{One has
$H^q(V,\mathcal O_V(-n))=0$ for all sufficiently large $n$ and
$0\leq q<p$.}

\phantomsection\label{fac:III:5:no:75:remark}
\textsc{Remark.} The
\hyperref[fac:III:5:no:75:thm:3:proof]{proof above} applies more generally whenever the
ideal $\mathcal I_x(V)$ admits a system of $r-p$ generators; in other
words, whenever $V$ is locally a complete intersection at every point.
